\section{Soft margin y regularización}

\begin{frame}{Soft margin: permitir violaciones con variables de holgura}
\footnotesize
\begin{columns}[T,onlytextwidth]
\begin{column}{0.56\textwidth}
\begin{ccdebox}
\centering
\textbf{Formulación primal usada por scikit-learn}
\[
\min_{\mathbf w,b,\boldsymbol\xi}
\quad
\frac{1}{2}\|\mathbf w\|^2
+ C\sum_{i=1}^{n}\xi_i
\]
\[
\text{s.a.}\quad
 y_i(\mathbf w^\top\mathbf x_i+b)\ge 1-\xi_i,
\qquad \xi_i\ge 0
\]
\end{ccdebox}
\vspace{0.12cm}
\begin{keybox}
$\xi_i$ mide cuánto puede violar el margen la observación $i$; si cruza la frontera, la violación es aún mayor.
\end{keybox}
\end{column}
\begin{column}{0.40\textwidth}
\begin{contrastbox}{ccdeaccent}{Dos objetivos en tensión}
\footnotesize
\begin{enumerate}
  \item Margen grande: minimizar $\|\mathbf w\|^2$.
  \item Pocas violaciones: minimizar $\sum_i \xi_i$.
\end{enumerate}
\end{contrastbox}
\vspace{0.14cm}
\begin{contrastbox}{ccdeblue}{El papel de $C$}
\footnotesize
$C$ decide cuánto penalizamos las violaciones frente a la anchura del margen.
\end{contrastbox}
\end{column}
\end{columns}
\source{Géron (2022), ``Soft Margin Classification'' y Eq. 5-2.}
\end{frame}

% ================================================================
\begin{frame}{$C$ controla el trade-off sesgo--varianza}
\small
\begin{columns}[T,onlytextwidth]
\begin{column}{0.47\textwidth}
\begin{contrastbox}{ccdeaccent}{En scikit-learn: $C$ pequeño}
\small
\begin{itemize}
  \item Más regularización.
  \item Margen más ancho.
  \item Tolera más violaciones.
  \item Menor varianza, más riesgo de underfitting.
\end{itemize}
\end{contrastbox}
\end{column}
\begin{column}{0.47\textwidth}
\begin{contrastbox}{warn}{En scikit-learn: $C$ grande}
\small
\begin{itemize}
  \item Menos regularización.
  \item Penaliza fuertemente violaciones.
  \item Frontera más ajustada a la muestra.
  \item Menor sesgo, más riesgo de overfitting.
\end{itemize}
\end{contrastbox}
\end{column}
\end{columns}
\vspace{0.22cm}
\begin{keybox}
\textbf{Cuidado con la notación:} ISLR (James et al.) usa otra parametrización de $C$ como presupuesto de violaciones, por lo que allí un $C$ mayor implica un margen más ancho. En scikit-learn y en Géron, $C$ es la penalización de las violaciones y la relación es la contraria.
\end{keybox}
\source{Géron (2022), cap. 5; James et al. (2013), sec. 9.2.2.}
\end{frame}

% ================================================================
\begin{frame}{Support vectors: los puntos que realmente fijan la frontera}
\small
\begin{columns}[c,onlytextwidth]
\begin{column}{0.50\textwidth}
\centering
\begin{tikzpicture}[scale=0.86]
  \draw[ccdelight,thick,dashed] (0.1,3.3)--(4.6,-0.2);
  \draw[ccdenavy,very thick] (0.65,4.05)--(5.15,0.55);
  \draw[ccdelight,thick,dashed] (1.2,4.8)--(5.7,1.3);
  \foreach \x/\y in {0.6/0.7,1.0/1.2,1.3/2.2,1.8/1.5,2.0/2.7}{\fill[ccdeaccent] (\x,\y) circle(2.6pt);}
  \foreach \x/\y in {3.5/2.1,4.0/3.0,4.4/2.7,4.7/3.6,5.0/4.1}{\fill[purplex] (\x,\y) circle(2.6pt);}
  \draw[warn,very thick] (2.02,2.68) circle(6pt);
  \draw[warn,very thick] (3.48,2.12) circle(6pt);
  \draw[warn,very thick] (1.34,2.19) circle(6pt);
  \node[text=warn,font=\scriptsize\bfseries] at (4.2,0.35){puntos activos};
  \draw[->,warn] (3.9,0.65)--(3.48,1.9);
\end{tikzpicture}
\end{column}
\begin{column}{0.46\textwidth}
\begin{itemize}
  \item Son observaciones sobre el margen, dentro del margen o mal clasificadas.
  \item En la solución dual tienen multiplicadores $\alpha_i>0$.
  \item Observaciones suficientemente alejadas del margen tienen $\alpha_i=0$.
  \item Por ello la función de decisión depende de un subconjunto de entrenamiento.
\end{itemize}
\vspace{0.14cm}
\begin{keybox}
Esta \textbf{esparsidad} es una propiedad importante: los puntos lejanos a la frontera no cambian la solución directamente.
\end{keybox}
\end{column}
\end{columns}
\source{James et al. (2013), secs. 9.1.3 y 9.2; Hastie et al. (2009), cap. 12.}
\end{frame}

% ================================================================
